% Transcription of Cayley's Collected Mathematical Papers, Volume IV
% PDF pages 376--400 (printed pages 354--378)
% Papers: 271 (conclusion), 272, 273, 274, 275 (start)

\documentclass[11pt]{article}
\usepackage[a4paper,margin=0.65in]{geometry}
\usepackage[T1]{fontenc}
\usepackage{amsmath,amssymb,amsfonts}
\usepackage{graphicx}
\usepackage{pdfpages}
\usepackage[english,shorthands=off]{babel}
\usepackage{fancyhdr}
\usepackage[bookmarks=false,colorlinks=false]{hyperref}
\usepackage[protrusion=true,expansion=false]{microtype}
\emergencystretch=3em
\tolerance=600
\providecommand{\modd}{\mathrm{mod}}
\providecommand{\Disc}{\mathrm{Disc}}
\providecommand{\canont}{}
\providecommand{\asterism}{*}
\providecommand{\Hessian}{H}
\providecommand{\tome}{}
\providecommand{\page}{p.}
\pagestyle{fancy}\fancyhf{}\fancyhead[L]{Pages 376--400}\fancyhead[R]{Cayley --- Collected Papers, Vol. IV}\renewcommand{\headrulewidth}{0.4pt}
\setlength{\headheight}{15pt}

\newcommand{\Cdel}{\mid}

\begin{document}

\thispagestyle{empty}
\begin{center}
\vspace*{2cm}
{\Large\scshape The Collected Mathematical Papers}\\[0.4em]
{\large of}\\[0.4em]
{\Large\scshape Arthur Cayley}\\[1.5em]
{\large Volume IV --- pages 354--378}\\[0.6em]
{\normalsize Paper 271 (conclusion), Papers 272, 273, 274, 275 (start)}\\[2em]
\vfill
\end{center}
\newpage

\setcounter{page}{354}

% ============================================================
% PAPER 271 (conclusion)
% SUR L'INVARIANT LE PLUS SIMPLE D'UNE FONCTION QUADRATIQUE BI-TERNAIRE
% ============================================================

\noindent\textbf{354}\hfill\textsc{sur l'invariant le plus simple d'une fonction \&c.}\hfill\textbf{[271}

\medskip

\noindent Le r\'eciprocant (multipli\'e par 2) est
\[
= -2 \begin{vmatrix}
\xi, & \eta, & \zeta \\
\xi, & \lambda, & \nu l, & \mu l \\
\eta, & \nu l, & \mu, & \lambda l \\
\zeta, & \mu l, & \lambda l, & \nu
\end{vmatrix},
\]
lequel, en forme d\'evelopp\'ee, donne le r\'esultat
\[
\begin{aligned}
&\phantom{+}(-2l^2 \nu \ldots \ldots \ldots + 2\mu \nu \ldots\ldots)\,\xi^2\\
&+(\ldots\,-2l^2 \nu \ldots \ldots \ldots + 2\nu\lambda\ldots)\,\eta^2\\
&+(\ldots \ldots -2l^2 \nu \ldots \ldots \ldots + 2\lambda\mu)\,\zeta^2\\
&+(-2l\lambda^2 \ldots \ldots \ldots \ldots + 2l\mu\nu\ldots)\,2\eta\zeta\\
&+(\ldots -2l\mu^2 \ldots \ldots +2l^2\nu\lambda \ldots)\,2\zeta\xi\\
&+(\ldots \ldots -2l\nu^2 \ldots \ldots +2l^2\lambda\mu)\,2\xi\eta,
\end{aligned}
\]
que l'on peut aussi exprimer comme fonction quadratique bi-ternaire, savoir:
\[
\begin{pmatrix}
-2l^2 & \cdot & \cdot & 1 & \cdot & \cdot \\
\cdot & -2l^2 & \cdot & \cdot & 1 & \cdot \\
\cdot & \cdot & -2l^2 & \cdot & \cdot & 1 \\
-2l & \cdot & \cdot & l^2 & \cdot & \cdot \\
\cdot & -2l & \cdot & \cdot & l^2 & \cdot \\
\cdot & \cdot & -2l & \cdot & \cdot & l^2
\end{pmatrix}(\lambda, \mu, \nu)^2 (\xi, \eta, \zeta)^2.
\]
En repr\'esentant l'invariant $2T$ par le d\'eveloppement incomplet donn\'e ci-dessus, on obtient dans le cas dont il s'agit
\[
\begin{aligned}
2T = {} & \begin{vmatrix}
-2l^2 & -2l^2 & -2l^2 \\
\cdot & -2l^2 & \cdot \\
\cdot & \cdot & -2l^2
\end{vmatrix}
- \begin{vmatrix}
-2l & \cdot & \cdot \\
\cdot & -2l & \cdot \\
\cdot & \cdot & -2l
\end{vmatrix}
+ \begin{vmatrix}
l^2 & \cdot & \cdot \\
\cdot & l^2 & \cdot \\
\cdot & \cdot & l^2
\end{vmatrix}
-\text{etc. etc.}\\
&{} + 2\begin{vmatrix}
-2l & \cdot & \cdot \\
\cdot & -2l & \cdot \\
\cdot & \cdot & -2l
\end{vmatrix},
\end{aligned}
\]

\newpage
\noindent\textbf{271]}\hfill\textsc{sur l'invariant le plus simple d'une fonction \&c.}\hfill\textbf{355}

\medskip

\noindent ce qui se r\'eduit \`a la valeur finale
\[
\begin{aligned}
& -8 l^6 + 2 \\
& \quad - (4 l^6 + 4 l^3 + 4 l^3) \\
& \quad - (4 l^3 + 4 l^3 + 4 l^3) \\
& \quad - (l^6 + l^6 + l^6) \\
& \quad + 2(-8 l^3 + 2 l^3) \\
& = 2(1 - 20 l^3 - 8 l^6).
\end{aligned}
\]
Les invariants $S$ et $T$ de la fonction cubique $x^3 + y^3 + z^3 + 6 l x y z$ sont
\[
\begin{aligned}
S &= -l + l^4,\\
T &= 1 - 20 l^3 - 8 l^6;
\end{aligned}
\]
l'invariant $2T$ de la fonction quadratique bi-ternaire est donc pr\'ecis\'ement \'egal \`a $2T$, o\`u $T$ est l'invariant qui vient d'\^etre donn\'e, et on a
\[
C_6 = 2 T.
\]
Le discriminant (multipli\'e par 6) de la fonction $\lambda U + \mu V + \nu W$ est
\[
= 6 \begin{vmatrix}
\lambda, & l\nu, & l\mu \\
l\nu, & \mu, & l\lambda \\
l\mu, & l\nu, & \nu
\end{vmatrix}
= 6\, [-l^2 \lambda^3 - l^2 \mu^3 - l^2 \nu^3 + (1 + 2 l^3)\, \lambda \mu \nu ],
\]
o\`u la fonction en parenth\`eses est ce que devient le Hessien de la fonction cubique $x^3 + y^3 + z^3 + 6 l x y z$, lorsqu'on y substitue $\lambda, \mu, \nu$ au lieu de $x, y, z$. Son invariant $S$ est
\[
\begin{aligned}
&= (1 + 2 l^3)^2\, 216 l^2 + (1 + 2 l^3)^4 \\
&= 1 + 8 l^3 + 240 l^6 + 464 l^9 + 16 l^{12}.
\end{aligned}
\]
Cette quantit\'e est en m\^eme temps un invariant de la fonction cubique
\[
x^3 + y^3 + z^3 + 6 l x y z,
\]
et comme telle elle doit s'exprimer en fonction des deux invariants $S$ et $T$ de cette derni\`ere: en effet elle se r\'eduit \`a
\[
(1 - 20 l^3 - 8 l^6)^2 - 48(-l + l^4)^3 = T^2 - 48 S^3.
\]
Nous avons donc
\[
C_{12} = T^2 - 48 S^3,
\]
et comme nous venons de trouver
\[
C_6 = 2 T,
\]
l'\'equation $12 R = 16 C_{12} - C_6^2$ se r\'eduit \`a
\[
\begin{aligned}
12 R &= 16(T^2 - 48 S^3) - 4 T^2 \\
&= 12 T^2 - 768 S^3,
\end{aligned}
\]

\hfill $45$--$2$

\newpage
\noindent\textbf{356}\hfill\textsc{sur l'invariant le plus simple d'une fonction \&c.}\hfill\textbf{[271}

\medskip

\noindent ou enfin \`a
\[
R = T^2 - 64 S^3,
\]
\'equation qui fait voir que le r\'esultant $R$ dont il s'agit est effectivement \'egal \`a $R = (1 + 8 l^3)^3$, c.~\`a~d. au discriminant de la fonction cubique
\[
x^3 + y^3 + z^3 + 6 l x y z,
\]
ce qui donne la v\'erification du th\'eor\`eme g\'en\'eral. Je m'\'etais servi d'abord d'une analyse moins simple, en consid\'erant le cas dans lequel on prend pour les trois fonctions les formes
\[
x^2 + y^2 + z^2 + 2 l y z + 2 m z x + 2 n x y.
\]

Il vaut la peine, je crois, de donner les expressions correspondantes de $C_6$, $C_{12}$, $R$. On a ici \`a consid\'erer la fonction
\[
\lambda(x^2 - y^2) + \mu(x^2 + y^2 + z^2 + 2 l y z + 2 m z x + 2 n x y) + \nu(z^2 - y^2),
\]
laquelle peut s'\'ecrire comme suit,
\[
(\mu + \lambda,\ \mu - \lambda - \nu,\ \mu + \nu,\ l\mu,\ m\mu,\ n\mu)(x, y, z)^2.
\]
Le r\'eciprocant (multipli\'e par 2) est
\[
= -2 \begin{vmatrix}
\xi, & \eta, & \zeta \\
\xi, & \mu+\lambda, & n\mu, & m\mu \\
\eta, & n\mu, & \mu-\lambda-\nu, & l\mu \\
\zeta, & m\mu, & l\mu, & \mu+\nu
\end{vmatrix}
\]
et en \'ecrivant $a$, $b$, $c$, $f$, $g$, $h$ au lieu de $1 - l^2$, $1 - m^2$, $1 - n^2$, $mn - l$, $nl - m$, $lm - n$, cette fonction est
\[
\begin{aligned}
= {}& (2 a \mu^2 - 2 \lambda \mu - 2 \nu \lambda - 2 l^2)\, \xi^2 \\
& + (2 b \mu^2 + 2 l^2 + 2 \mu^2 + 2 \nu \lambda)\, \eta^2 \\
& + (2 c \mu^2 - 2 \mu \lambda - 2 \nu \lambda - 2 \lambda^2)\, \zeta^2 \\
& + (2 f \mu^2 - 2 l \lambda \mu)\, 2 \eta \zeta \\
& + (2 g \mu^2 + 2 m \lambda \mu + 2 m l \nu \mu)\, 2 \zeta \xi \\
& + (2 h \mu^2 - 2 n \nu \mu)\, 2 \xi \eta,
\end{aligned}
\]
ou, \'ecrite en forme de fonction quadratique bi-ternaire,
\[
=\begin{pmatrix}
0, & 2a, & -2, & 0, & 1, & -1 \\
0, & 2b, & 0, & 1, & 1, & 1 \\
-2, & 2c, & 0, & -1, & -1, & 0 \\
0, & 2f, & 0, & 0, & 0, & -l \\
0, & 2g, & 0, & m, & 0, & l \\
0, & 2h, & 0, & -n, & 0, & 0
\end{pmatrix}(\lambda, \mu, \nu)^2 (\xi, \eta, \zeta)^2.
\]

\newpage
\noindent\textbf{271]}\hfill\textsc{sur l'invariant le plus simple d'une fonction \&c.}\hfill\textbf{357}

\medskip

\noindent En exprimant comme auparavant l'invariant $2T$ par une somme de cinq termes, on trouve sans peine que ces termes sont respectivement
\[
18 - 4 l^2 - 4 m^2 - 4 n^2,\quad -4 l^2,\ -4 m^2,\ -4 n^2,\ 0;
\]
l'invariant sera donc $-8(l^2 + m^2 + n^2) + 18$; et l'on a pour le combinant du sixi\`eme ordre
\[
C_6 = -8(l^2 + m^2 + n^2) + 18.
\]
Le discriminant (multipli\'e par 6) de la fonction $\lambda U + \mu V + \nu W$ est
\[
= 6 \begin{vmatrix}
\mu+\lambda, & n\mu, & m\mu \\
n\mu, & \mu-\lambda-\nu, & l\mu \\
m\mu, & l\mu, & \mu+\nu
\end{vmatrix}
\]
\[
\begin{aligned}
={}& 6\,(1 - l^2 - m^2 - n^2 + 2 l m n)\, \mu^3 + 0 \cdot \nu^3 \\
&+ 6\,(m^2 - n^2)\, \mu^2 \nu - 6 \lambda^2 \mu - 6 \lambda \nu^2 - 6 \mu^2 \lambda - 6 \nu \lambda^2 + 6\,(m^2 - l^2)\, \lambda \mu^2 - 6 \lambda \mu \nu \\
={}& \big(0,\ 6(1 - l^2 - m^2 - n^2 + 2 l m n),\ 0,\ 2(m^2 - n^2),\ -2,\ -2,\ -2,\ -2,\ 2(m^2 - l^2),\ -1\big)(\lambda, \mu, \nu)^3 \\
={}& (a, b, c, i, j, k, l_1, j_1, h, 1)(\lambda, \mu, \nu)^3,
\end{aligned}
\]
o\`u l'on a pos\'e $a = 0$, $b = 6(1 - l^2 - m^2 - n^2 + 2 l m n)$, etc. L'invariant $S$ de cette derni\`ere fonction est
\[
\begin{aligned}
={}& 1 \\
&\quad - 2\,(-4(m^2 - n^2) + 4 - 4(m^2 - n^2)) \\
&\quad - 3\,(8(m^2 - n^2) + 8(m^2 - l^2)) \\
&\quad - 4 b \\
&\quad + 16\,(m^2 - n^2) - 16\,(m^2 - n^2)(m^2 - l^2) + 16\,(m^2 - l^2) \\
&\quad + 16\,(m^2 - n^2)^2 + 16 + 16\,(m^2 - l^2) \\
&\quad + 16 b \\
={}& 9 + 12 b + 16(l^4 + m^4 + n^4 - m^2 n^2 - n^2 l^2 - l^2 m^2),
\end{aligned}
\]
ou, si l'on substitue la valeur de $b$,
\[
C_{12} = 16\,(l^4 + m^4 + n^4 - m^2 n^2 - n^2 l^2 - l^2 m^2) + 144\, l m n - 72\,(l^2 + m^2 + n^2) + 81.
\]
On vient de trouver
\[
C_6 = -8(l^2 + m^2 + n^2) + 18,
\]
\'equation qui donne
\[
C_6{}^2 = 64\,(l^4 + m^4 + n^4 + 2 m^2 n^2 + 2 n^2 l^2 + 2 l^2 m^2) - 288\,(l^2 + m^2 + n^2) + 324.
\]

\newpage
\noindent\textbf{358}\hfill\textsc{sur l'invariant le plus simple d'une fonction \&c.}\hfill\textbf{[271}

\medskip

\noindent De ces r\'esultats combin\'es on conclut
\[
16 C_{12} - C_6{}^2 = 192\,(l^4 + m^4 + n^4 - 2 m^2 n^2 - 2 n^2 l^2 - 2 l^2 m^2) + 2304\, l m n - 864\,(l^2 + m^2 + n^2) + 972.
\]
\'Evidemment l'expression du r\'esultant est, \`a un facteur num\'erique pr\`es,
\[
R = (3 + 2 l + 2 m + 2 n)(3 + 2 l - 2 m - 2 n)(3 - 2 l + 2 m - 2 n)(3 - 2 l - 2 m + 2 n),
\]
ou en r\'eduisant
\[
R = 16\,(l^4 + m^4 + n^4 - 2 m^2 n^2 - 2 n^2 l^2 - 2 l^2 m^2) + 192\, l m n - 72\,(l^2 + m^2 + n^2) + 81.
\]
On retrouve donc
\[
12 R = 16 C_{12} - C_6{}^2,
\]
\'equation qu'il s'agissait de v\'erifier.

\medskip

En terminant je remarque que la plus grande partie de ce m\'emoire est tir\'ee d'un manuscrit dat\'e ``Machynlleth, 13 Ao\^ut 1853.''

\medskip

\hfill Londres, 22 Mars 1859.

% ============================================================
% PAPER 272
% DEMONSTRATION D'UN THEOREME DE JACOBI PAR RAPPORT AU PROBLEME DE PFAFF
% ============================================================

\newpage
\noindent\textbf{272]}\hfill\textbf{359}

\medskip

\begin{center}
{\Large\bfseries 272.}\\[1em]
{\large\scshape D\'emonstration d'un th\'eor\`eme de Jacobi par rapport au probl\`eme de Pfaff.}\\[1em]
{\small [From the \textit{Journal f\"ur die reine und angewandte Mathematik} (Crelle), tom. \textsc{lvii}.\ (1860), pp.~273--277.]}
\end{center}

\medskip

\noindent Dans le m\'emoire de Jacobi ``Theoria novi multiplicatoris etc.'' (t.~\textsc{xxix}.\ de ce Journal, 1845) on trouve p.~253 le passage que voici: ``Methodum ad solvendum problema Pfaffianum ab ipso auctore adhibitam, data occasione observo per plures et altiores procedere integrationes quam methodus vera et genuina poscat. Quam novam methodum pro exemplo simplice explicabo. Ad aequationem differentialem
\[
X_1 dx_1 + X_2 dx_2 + X_3 dx_3 + X_4 dx_4 = 0
\]
per duas aequationes integrandam poscit Pfaffiana methodus\ldots\ integrationem completam systematis trium aequationum differentialium primi ordinis inter quatuor variabiles ac deinde unius aequationis differentialis primi ordinis inter duas variabiles. Illius igitur systematis integrali uno invento, secundum illam methodum restat integratio completa duarum aequationum differentialium primi ordinis inter tres variabiles sive unius aequationis differentialis secundi ordinis inter duas variabiles ac deinde aequationis differentialis primi ordinis inter duas variabiles. At observo si integrali illo invento exprimatur $x_4$ per $x_1, x_2, x_3$, aequationem differentialem propositam abire in aliam linearem primi ordinis inter tres variabiles conditioni integrabilitatis satisfacientem; cujus integrationem vidimus (voir p.~246) absolvi posse per integrationes separatas duarum aequationum differentialium primi ordinis inter duas variabiles. Unde in locum aequationis differentialis secundi ordinis, tantum integrandae sunt duae aequationes differentiales separatae primi ordinis, quae est reductio maxime insignis; integrationi autem aequationis differentialis primi ordinis postremo praestandae omnino supersedetur. Tractatio hujus rei gravissimae completa ac generalis alii commentationi reservanda

\newpage
\noindent\textbf{360}\hfill\textsc{d\'emonstration d'un th\'eor\`eme de jacobi}\hfill\textbf{[272}

\medskip

\noindent est.'' Le th\'eor\`eme dont il s'agit est \'enonc\'e dans les mots ``conditioni integrabilitatis satisfacientem,'' mais j'ai cit\'e le passage en entier pour montrer l'importance qu'\`a bon droit Jacobi attachait \`a la th\'eorie dont ce th\'eor\`eme fait partie. On sait que le probl\`eme de la solution d'une \'equation \`a diff\'erences partielles du premier ordre n'est qu'un cas particulier du probl\`eme de Pfaff et que la solution des \'equations diff\'erentielles de la Dynamique et celle d'autres syst\`emes est intimement li\'ee avec la question de la solution de ce cas particulier: le probl\`eme de Pfaff est ainsi de la plus haute importance dans l'analyse. Jacobi a parl\'e autre part d'un m\'emoire sur la m\'ecanique analytique dont il s'est beaucoup occup\'e, lequel aurait naturellement des rapports avec celui-l\`a; malheureusement ni l'un ni l'autre de ces m\'emoires n'ont paru.

\medskip

Pour d\'emontrer le th\'eor\`eme, consid\'erons l'expression diff\'erentielle
\[
\nabla = X_1\, dx_1 + X_2\, dx_2 + X_3\, dx_3 + X_4\, dx_4.
\]
En supposant que $a = \phi(x_1, x_2, x_3, x_4)$ soit une int\'egrale des \'equations diff\'erentielles de Pfaff, si au moyen de cette \'equation on exprime $x_4$ en fonction de $(x_1, x_2, x_3, a)$ l'expression $\nabla$ prendra la forme
\[
\nabla_1 = Y_1\, dx_1 + Y_2\, dx_2 + Y_3\, dx_3 + A\, da
\]
et il s'agit de faire voir que l'\'equation
\[
Y_1\, dx_1 + Y_2\, dx_2 + Y_3\, dx_3 = 0
\]
\`a laquelle $\nabla_1 = 0$ se r\'eduit, lorsqu'on fait $a = $ constante, sera int\'egrable au moyen d'un facteur, ou autrement dit, que l'expression $Y_1\, dx_1 + Y_2\, dx_2 + Y_3\, dx_3$ sera de la forme $U du$, c.~\`a~d. qu'il existe une quantit\'e $u$, fonction de $x_1, x_2, x_3, a$, telle qu'en traitant $a$ comme une constante on a
\[
Y_1\, dx_1 + Y_2\, dx_2 + Y_3\, dx_3 = U du.
\]
Ainsi dans cette formule $du$ denote $\dfrac{du}{dx_1}\, dx_1 + \dfrac{du}{dx_2}\, dx_2 + \dfrac{du}{dx_3}\, dx_3$, et si l'on veut regarder $a$ comme variable il faudrait y \'ecrire $du - \dfrac{du}{da}\, da$ au lieu de $du$; on aurait ainsi
\[
\nabla_1 = U\!\left(du - \frac{du}{da}\, da\right) + A\, da,
\]
ou en changeant la signification de $A$,
\[
\nabla = X_1\, dx_1 + X_2\, dx_2 + X_3\, dx_3 + X_4\, dx_4 = U du + A da,
\]
o\`u $u, a$, et de m\^eme $U, A$, seront des fonctions donn\'ees de $x_1, x_2, x_3, x_4$; c'est en effet la forme \`a laquelle dans le probl\`eme de Pfaff il s'agit de r\'eduire l'expression diff\'erentielle $\nabla$.

\medskip

Or les \'equations diff\'erentielles de Pfaff sont le syst\`eme que voici, savoir en \'ecrivant
\[
12 = \frac{d X_2}{d x_1} - \frac{d X_1}{d x_2},\ \text{etc.},
\]

\newpage
\noindent\textbf{272]}\hfill\textsc{par rapport au probl\`eme de pfaff.}\hfill\textbf{361}

\medskip

\noindent ce qui implique $12 = -21$, $11 = 0$, etc., et en introduisant la variable auxiliaire $t$, qui doit \^etre \'elimin\'ee, le syst\`eme est
\[
\begin{aligned}
0 &= X_1\, dx_1 + X_2\, dx_2 + X_3\, dx_3 + X_4\, dx_4,\\
X_1\, dt &= \phantom{*}\ \ {} * \ \ 12\, dx_2 + 13\, dx_3 + 14\, dx_4,\\
X_2\, dt &= 21\, dx_1 \ {} * \ {} + 23\, dx_3 + 24\, dx_4,\\
X_3\, dt &= 31\, dx_1 + 32\, dx_2 \ \ {} * \ {} + 34\, dx_4,\\
X_4\, dt &= 41\, dx_1 + 42\, dx_2 + 43\, dx_3 \ \ {} *,
\end{aligned}
\]
lequel ne contient que quatre \'equations ind\'ependantes. On donne \`a ce syst\`eme une forme plus sym\'etrique en \'ecrivant $x_0$ au lieu de $t$, et en y mettant de plus
\[
X_1 = 01 = -10,\ \text{etc.}
\]
les \'equations deviennent par ce moyen
\[
\begin{aligned}
0 &= \phantom{*}\ \ 01\, dx_1 + 02\, dx_2 + 03\, dx_3 + 04\, dx_4,\\
0 &= 10\, dx_0 \ \ {} * \ {} + 12\, dx_2 + 13\, dx_3 + 14\, dx_4,\\
0 &= 20\, dx_0 + 21\, dx_1 \ {} * \ {} + 23\, dx_3 + 24\, dx_4,\\
0 &= 30\, dx_0 + 31\, dx_1 + 32\, dx_2 \ {} * \ {} + 34\, dx_4,\\
0 &= 40\, dx_0 + 41\, dx_1 + 42\, dx_2 + 43\, dx_3 \ \ {} *,
\end{aligned}
\]
et on d\'eduit de l\`a
\[
dx_0 : dx_1 : dx_2 : dx_3 : dx_4 = 1234. : 2340. : 3401. : 4012. : 0123.,
\]
o\`u $1234.$, etc.\ sont les fonctions de Jacobi que j'ai nomm\'ees ``Pfaffiens,'' savoir on a
\[
\begin{aligned}
12.34 &= 12 \cdot 34 + 13 \cdot 42 + 14 \cdot 23,\\
2340. &= 23 \cdot 40 + 24 \cdot 03 + 20 \cdot 34, \quad = \ \ {} * - X_2 34 - X_3 42 - X_4 23,\\
3401. &= 34 \cdot 01 + 30 \cdot 14 + 31 \cdot 40, \quad = X_1 34 \ {} * + X_3 41 + X_4 13,\\
4012. &= 40 \cdot 12 + 41 \cdot 20 + 42 \cdot 01, \quad = -X_1 24 - X_2 41 \ {} * - X_4 12,\\
0123. &= 01 \cdot 23 + 02 \cdot 31 + 03 \cdot 12, \quad = X_1 23 + X_2 31 + X_3 12 \ {} * .
\end{aligned}
\]
Cette transformation se trouve en effet dans le m\'emoire de Jacobi ``Ueber die Pfaffsche Methode etc.'' (t.~\textsc{ii}.\ de ce Journal p.~347 et suivantes, 1827), o\`u cependant Jacobi ne s'est pas servi de l'algorithme $X_1 = 01 = -10$ etc.\ au moyen duquel la forme de la solution devient si simple.

\medskip

Faisons attention \`a pr\'esent \`a ce que l'\'equation
\[
a = \phi(x_1, x_2, x_3, x_4)
\]
est suppos\'ee \^etre une int\'egrale des \'equations diff\'erentielles. On en d\'eduit d'abord $x_4 = \text{fonc.}(x_1, x_2, x_3, a)$, et de l\`a en traitant $a$ comme constante
\[
dx_4 = \frac{d x_4}{d x_1}\, dx_1 + \frac{d x_4}{d x_2}\, dx_2 + \frac{d x_4}{d x_3}\, dx_3,
\]

\hfill C, IV.\hfill 46

\newpage
\noindent\textbf{362}\hfill\textsc{d\'emonstration d'un th\'eor\`eme de jacobi}\hfill\textbf{[272}

\medskip

\noindent \'equation qui doit \^etre satisfaite par les valeurs qu'on vient de trouver pour $dx_0 : dx_1 : dx_2 : dx_3 : dx_4$, et qui par cons\'equent donne:
\[
0 = 0123. - \frac{dx_4}{dx_1}\, 2340. - \frac{dx_4}{dx_2}\, 3401. - \frac{dx_4}{dx_3}\, 4012.
\]
En y substituant pour $0123.$, etc.\ leurs valeurs, modifi\'ees en sorte que tous les termes deviennent positifs (ce qui se fait \`a l'aide des relations $12 = -21$ etc.), on trouve
\[
\begin{aligned}
& \frac{dx_4}{dx_1}\,(\ \ \ * \ \ X_2 34 + X_3 42 + X_4 23) \\
& + \frac{dx_4}{dx_2}\,(X_1 43 \ \ \ * \ \ + X_3 14 + X_4 31) \\
& + \frac{dx_4}{dx_3}\,(X_1 24 + X_2 41 \ \ \ * \ \ + X_4 12) \\
& \quad + (X_1 32 + X_2 13 + X_3 21\ \ \ * \ \ ) = 0,
\end{aligned}
\]
\'equation aux diff\'erences partielles \`a laquelle satisfait la variable $x_4$, lorsqu'au moyen de l'\'equation $a = \phi(x_1, x_2, x_3, x_4)$ (qui est une int\'egrale des \'equations diff\'erentielles de Pfaff) elle est donn\'ee en fonction de $x_1, x_2, x_3$. Nous allons voir que c'est cette derni\`ere \'equation, dans laquelle est contenu le th\'eor\`eme dont la d\'emonstration fait l'objet de cette note. En effet l'expression diff\'erentielle
\[
\nabla = X_1\, dx_1 + X_2\, dx_2 + X_3\, dx_3 + X_4\, dx_4
\]
ayant \'et\'e transform\'ee en
\[
\nabla_1 = Y_1\, dx_1 + Y_2\, dx_2 + Y_3\, dx_3 + A\, da
\]
par la substitution de $a$ au lieu de $x_4$, il s'en suit qu'on a les relations:
\[
\begin{aligned}
Y_1 &= X_1 + X_4 \frac{dx_4}{dx_1},\\
Y_2 &= X_2 + X_4 \frac{dx_4}{dx_2},\\
Y_3 &= X_3 + X_4 \frac{dx_4}{dx_3},\\
A &= X_4 \frac{dx_4}{da},
\end{aligned}
\]
la variable $x_4$ contenue dans les quantit\'es $X_1, X_2, X_3, X_4$ doit \^etre remplac\'ee par son expression en fonction de $x_1, x_2, x_3, a$.

\medskip

Cela pos\'e, la condition pour que l'\'equation
\[
Y_1\, dx_1 + Y_2\, dx_2 + Y_3\, dx_3 = 0
\]
soit int\'egrable par un facteur \'etant d\'esign\'ee par
\[
Z = 0,
\]

\newpage
\noindent\textbf{272]}\hfill\textsc{par rapport au probl\`eme de pfaff.}\hfill\textbf{363}

\medskip

\noindent on sait que
\[
Z = Y_1\!\left(\frac{dY_2}{dx_3} - \frac{dY_3}{dx_2}\right) + Y_2\!\left(\frac{dY_3}{dx_1} - \frac{dY_1}{dx_3}\right) + Y_3\!\left(\frac{dY_1}{dx_2} - \frac{dY_2}{dx_1}\right).
\]
Mais d'apr\`es les relations que l'on vient d'\'etablir on a
\[
\frac{dY_2}{dx_3} = \frac{dX_2}{dx_3} + \frac{dX_2}{dx_4}\, \frac{dx_4}{dx_3} + \!\left(\frac{dX_4}{dx_3} + \frac{dX_4}{dx_4}\, \frac{dx_4}{dx_3}\right)\!\frac{dx_4}{dx_2} + X_4\, \frac{d^2 x_4}{dx_2\, dx_3},
\]
\[
\frac{dY_3}{dx_2} = \frac{dX_3}{dx_2} + \frac{dX_3}{dx_4}\, \frac{dx_4}{dx_2} + \!\left(\frac{dX_4}{dx_2} + \frac{dX_4}{dx_4}\, \frac{dx_4}{dx_2}\right)\!\frac{dx_4}{dx_3} + X_4\, \frac{d^2 x_4}{dx_2\, dx_3},
\]
d'o\`u l'on tire
\[
\frac{dY_2}{dx_3} - \frac{dY_3}{dx_2} = 23 + 24\, \frac{dx_4}{dx_3} + 43\, \frac{dx_4}{dx_2},
\]
et on trouve de m\^eme
\[
\frac{dY_3}{dx_1} - \frac{dY_1}{dx_3} = 31 + 34\, \frac{dx_4}{dx_1} + 41\, \frac{dx_4}{dx_3},
\]
\[
\frac{dY_1}{dx_2} - \frac{dY_2}{dx_1} = 12 + 14\, \frac{dx_4}{dx_2} + 42\, \frac{dx_4}{dx_1};
\]
donc en ajoutant ces \'equations apr\`es les avoir multipli\'ees par $X_1 + X_4 \frac{dx_4}{dx_1}$, $X_2 + X_4 \frac{dx_4}{dx_2}$, $X_3 + X_4 \frac{dx_4}{dx_3}$, les termes qui contiennent les produits $\frac{dx_4}{dx_1}\, \frac{dx_4}{dx_2}$, etc.\ se d\'etruisent, et l'on obtient l'\'equation finale
\[
\begin{aligned}
Z = {}& \frac{dx_4}{dx_1}\,(\ \ * \ \ X_2 34 + X_3 42 + X_4 23) \\
& + \frac{dx_4}{dx_2}\,(X_1 43 \ \ * \ \ + X_3 14 + X_4 31) \\
& + \frac{dx_4}{dx_3}\,(X_1 24 + X_2 41 \ \ * \ \ + X_4 12) \\
& \quad + (X_1 32 + X_2 13 + X_3 21 \ \ *),
\end{aligned}
\]
dont la seconde partie s'\'evanouit comme il a \'et\'e d\'emontr\'e ci-dessus. Par cons\'equent la condition d'int\'egrabilit\'e $Z = 0$ se trouve remplie, ce qu'il s'agissait de prouver.

\medskip

\hfill Londres, 3 Septembre 1859.\hfill $46$--$2$

% ============================================================
% PAPER 273
% NOTE SUR LA TRANSFORMATION DE TSCHIRNHAUSEN
% ============================================================

\newpage
\noindent\textbf{364}\hfill\textbf{[273}

\medskip

\begin{center}
{\Large\bfseries 273.}\\[1em]
{\large\scshape Note sur la transformation de Tschirnhausen.}\\[1em]
{\small [From the \textit{Journal f\"ur die reine und angewandte Mathematik} (Crelle), tom. \textsc{lviii}.\ (1861), pp.~259--262.]}
\end{center}

\medskip

\noindent On trouve dans le m\'emoire de M.~Hermite ``Sur quelques th\'eor\`emes d'alg\`ebre et la r\'esolution de l'\'equation du quatri\`eme degr\'e'' (\textit{Comptes Rendus}, t.~\textsc{xlvi}.\ p.~961, 1858) un th\'eor\`eme tr\`es-important relatif \`a la transformation de Tschirnhausen, \`a l'aide de laquelle une \'equation $f(x) = 0$ est transform\'ee en une autre du m\^eme degr\'e en $y$ par la substitution $y = \phi x$, o\`u $\phi x$ d\'esigne une fonction rationnelle. En consid\'erant pour fixer les id\'ees une \'equation du quatri\`eme degr\'e
\[
(a, b, c, d, e \Cdel x, 1)^4 = 0,
\]
M.~Hermite donne \`a l'\'equation $y = \phi x$ la forme que voici,
\[
y = a T + (a x + 4 b)\, T_0 + (a x^2 + 4 b x + 6 c)\, T_1 + (a x^3 + 4 b x^2 + 6 c x + 4 d)\, T_2,
\]
et il fait voir que les coefficients de la transform\'ee en $y$ ont la propri\'et\'e suivante: toute fonction de ces coefficients, laquelle exprim\'ee en fonction de $a, b, c, d, e, T, T_0, T_1, T_2$ ne contient pas $T$, est un invariant, c.~\`a~d.\ un invariant des deux fonctions
\[
(a, b, c, d, e \Cdel \xi, \eta)^4,\quad (T_0, T_1, T_2 \Cdel \eta, -\xi)^2.
\]
Cela revient \`a dire qu'en supposant la valeur de $T$ d\'etermin\'ee de mani\`ere \`a faire \'evanouir dans l'\'equation en $y$ le coefficient du second terme (de $y^3$), les autres coefficients seront des invariants, de sorte que, si dans le polynome en $y$ qui est \'egal \`a z\'ero on consid\`ere $y$ comme une constante absolue, le polynome tout entier sera un invariant des deux fonctions ci-dessus mentionn\'ees. On trouve sans peine la valeur que doit avoir $T$, elle est donn\'ee par l'\'equation
\[
0 = a T + 3 b T_0 + 3 c T_1 + d T_2,
\]

\newpage
\noindent\textbf{273]}\hfill\textsc{note sur la transformation de tschirnhausen.}\hfill\textbf{365}

\medskip

\noindent ce qui conduit pour $y$ \`a la valeur
\[
y = (a x + b)\, T_0 + (a x^2 + 4 b x + 3 c)\, T_1 + (a x^3 + 4 b x^2 + 6 c x + 3 d)\, T_2,
\]
et en m\^eme temps la transform\'ee en $y$ aura la propri\'et\'e dont il s'agit.

\medskip

Par rapport \`a la forme de l'expression que l'on vient de trouver pour $y$ il est bon de remarquer que les coefficients num\'eriques qu'on y rencontre, hormis ceux du terme en $x^0$, ou de $b T_0 + 3 c T_1 + 3 d T_2$, sont les coefficients de la puissance $(1, 1)^4$, tandis que les coefficients qui ont \'et\'e d\'esign\'es comme faisant exception sont ceux de la puissance $(1, 1)^3$. Une remarque pareille s'applique au cas g\'en\'eral. Pour d\'emontrer le th\'eor\`eme \'enonc\'e, je repr\'esente l'\'equation qui vient d'\^etre \'ecrite par $y = V$, la transform\'ee en $y$ sera donc
\[
(y - V_1)(y - V_2)(y - V_3)(y - V_4) = 0,
\]
o\`u $V_1, V_2, V_3, V_4$ sont ce que devient $V$ lorsqu'on y substitue successivement pour $x$ les quatre racines de l'\'equation $(a, b, c, d, e \Cdel x, 1)^4 = 0$. Or, en consid\'erant $y$ comme une constante, pour que l'expression qui forme la premi\`ere partie de l'\'equation que l'on vient d'\'ecrire soit un invariant, les conditions \`a remplir consistent en ce que cette expression se r\'eduise \`a z\'ero au moyen de l'un et l'autre des op\'erateurs
\[
a \partial_b + 2 b \partial_c + 3 c \partial_d + 4 d \partial_e - (T_1 \partial_{T_0} + 2 T_2 \partial_{T_1}),
\]
\[
4 b \partial_a + 3 c \partial_b + 2 d \partial_c + e \partial_d - (2 T_0 \partial_{T_1} + T_1 \partial_{T_2}).
\]
Ces conditions seront satisfaites si chacun des facteurs $y - V_1$, etc.\ est dou\'e de cette m\^eme propri\'et\'e, c.~\`a~d.\ si $y - V$, ou plus simplement si $V$, en y faisant $x$ \'egal \`a l'une des racines de l'\'equation en $x$, se r\'eduit \`a z\'ero au moyen de l'un et l'autre des deux op\'erateurs ci-dessus \'ecrits. Je consid\`ere le premier des deux op\'erateurs et pour abr\'eger je le d\'esigne par
\[
\Delta - (T_1 \partial_{T_0} + 2 T_2 \partial_{T_1}).
\]
Pour avoir $\Delta V$, il faut d'abord former la valeur de $\Delta x$. On l'obtient en op\'erant avec $\Delta$ sur l'\'equation $(a, b, c, d, e \Cdel x, 1)^4 = 0$, ce qui donne
\[
(a, b, c, d \Cdel x, 1)^3\, \Delta x + (a, b, c, d \Cdel x, 1)^3 = 0,\quad \text{ou } \Delta x = -1.
\]
La partie de $\Delta V$ qui d\'epend de la variation de $x$ est par cons\'equent
\[
-a T_1 + (-2 a x - 4 b)\, T_1 + (-3 a x^2 - 8 b x - 6 c)\, T_2.
\]
Pour l'autre partie de $\Delta V$ on trouve ais\'ement
\[
a T_0 + (4 a x + 6 b)\, T_1 + (4 a x^2 + 12 b x + 9 c)\, T_2,
\]
et de l\`a en ajoutant
\[
\Delta V = 2(a x + b)\, T_1 + (a x^2 + 4 b x + 3 c)\, T_2,
\]
ce qui est pr\'ecis\'ement \'egal \`a
\[
(T_1 \partial_{T_0} + 2 T_2 \partial_{T_1})\, V.
\]

\newpage
\noindent\textbf{366}\hfill\textsc{note sur la transformation de tschirnhausen.}\hfill\textbf{[273}

\medskip

\noindent Donc $V$ se r\'eduit \`a z\'ero par l'op\'erateur
\[
\Delta - (T_1 \partial_{T_0} + 2 T_2 \partial_{T_1}).
\]
De m\^eme en repr\'esentant le second op\'erateur par
\[
\nabla - (2 T_0 \partial_{T_1} + T_1 \partial_{T_2})
\]
on trouve d'abord
\[
(a, b, c, d \Cdel x, 1)^3\, \nabla x + x\,(b, c, d, e \Cdel x, 1)^3 = 0,
\]
mais en ayant \'egard \`a l'\'equation $(a, b, c, d, e \Cdel x, 1)^4 = 0$ la valeur de $\nabla x$ se r\'eduit \`a $\nabla x = x^2$. La partie de $\nabla V$ due \`a la variation de $x$ est par cons\'equent
\[
a x^2 T_0 + (2 a x^3 + 4 b x^2)\, T_1 + (3 a x^4 + 8 b x^3 + 6 c x^2)\, T_2.
\]
L'autre partie de $\nabla V$ est
\[
(4 b x + 3 c)\, T_0 + (4 b x^2 + 12 c x + 6 d)\, T_1 + (4 b x^3 + 12 c x^2 + 12 d x + 3 e)\, T_2.
\]
En les ajoutant, le coefficient de $T_2$ s'\'evanouit en vertu de l'\'equation en $x$, et l'on trouve
\[
\nabla V = (a x^2 + 4 b x + 3 c)\, T_0 + 2(a x^3 + 4 b x^2 + 6 c x + 3 d)\, T_1,
\]
ce qui est pr\'ecis\'ement \'egal \`a
\[
(2 T_0 \partial_{T_1} + T_1 \partial_{T_2})\, V.
\]
$V$ se r\'eduit donc \`a z\'ero au moyen de l'op\'erateur
\[
\nabla - (2 T_0 \partial_{T_1} + T_1 \partial_{T_2}),
\]
ce qui ach\`eve la d\'emonstration dont il s'agissait. Il va sans dire que la d\'emonstration serait conduite d'une mani\`ere semblable pour une \'equation de degr\'e quelconque. Pour avoir l'exemple le plus simple, je prends les \'equations
\[
(a, b, c, d \Cdel x, 1)^3 = 0,
\]
\[
y = (a x + b)\, T_0 + (a x^2 + 3 b x + 2 c)\, T_1,
\]
et pour effectuer l'\'elimination j'\'ecris
\[
y x = (a x^2 + b x)\, T_0 + (-c x - d)\, T_1,
\]
\[
y x^2 = (-2 b x^2 - 3 c x - d)\, T_0 + (-c x^2 - d x)\, T_1.
\]
Maintenant on a les trois \'equations
\[
\begin{aligned}
0 &= b T_0 + 2 c T_1 - y + x\,(a T_0 + 3 b T_1) + x^2 \cdot a T_0,\\
0 &= -d T_1 + x\,(b T_0 - c T_1 - y) + x^2 \cdot a T_0,\\
&\phantom{=}\quad + x\,(-3 c T_0 - d T_1)\\
0 &= -d T_0 \quad + x^2\,(-2 b T_0 - c T_1 - y),
\end{aligned}
\]
donc l'\'equation en $y$ est
\[
\begin{vmatrix}
y - b T_0 - 2 c T_1, & -a T_0 - 3 b T_1, & -a T_0 \\
d T_1, & y - b T_0 + c T_1, & -a T_0 \\
d T_0, & 3 c T_0 + d T_1, & y + 2 b T_0 + c T_1
\end{vmatrix} = 0.
\]

\newpage
\noindent\textbf{273]}\hfill\textsc{note sur la transformation de tschirnhausen.}\hfill\textbf{367}

\medskip

\noindent En ordonnant ce d\'eterminant suivant les puissances de $y$, les coefficients de $y^3$, $y^2$, $y$, $y^0$ deviennent des formes binaires en $T_0, T_1$ des ordres 0, 1, 2, 3. En calculant les valeurs de ces quatre coefficients on les trouve respectivement
\[
\begin{aligned}
&= 1,\\
&= 0,\\
&= 3\,(ac - b^2,\ ad - bc,\ bd - c^2 \Cdel T_0, T_1)^2,\\
&= (a^2 d - 3 a b c + 2 b^3,\ 3 a b d - 6 a c^2 + 3 b^2 c,\ -3 a c d + 6 b^2 d - 3 b c^2,\ -a d^2 + 3 b c d - 2 c^3 \Cdel T_0, T_1)^3,
\end{aligned}
\]
c'est-\`a-dire que l'\'equation en $y$ est celle-ci:
\[
\resizebox{\textwidth}{!}{$\displaystyle
y^3 + 3 y\,(a c - b^2,\ a d - b c,\ b d - c^2 \Cdel T_0, T_1)^2 + (a^2 d - 3 a b c + 2 b^3,\ 3 a b d - 6 a c^2 + 3 b^2 c,\ -3 a c d + 6 b^2 d - 3 b c^2,\ -a d^2 + 3 b c d - 2 c^3 \Cdel T_0, T_1)^3 = 0,
$}
\]
\'equation qui remplit en effet la condition ci-dessus mentionn\'ee, d'avoir pour coefficients des invariants des deux formes:
\[
(a, b, c, d \Cdel \xi, \eta)^3,\ (T_0, T_1 \Cdel \eta, -\xi).
\]
Dans le cas particulier dont il s'agit, la fonction $(T_0, T_1 \Cdel \eta, -\xi)$ est lin\'eaire, et l'on peut m\^eme dire que les coefficients sont des covariants de la seule fonction $(a, b, c, d \Cdel T_0, T_1)^3$ en y consid\'erant $T_0, T_1$ comme les variables.

\medskip

J'ai cru qu'il y avait de l'int\'er\^et de donner cette v\'erification. D'ailleurs je remarque qu'au moyen du th\'eor\`eme m\^eme on aurait pu trouver tout de suite l'\'equation en $y$, en \'ecrivant d'abord $T_1 = 0$, ce qui donne le syst\`eme
\[
(a, b, c, d \Cdel x, 1)^3 = 0,
\]
\[
y = (a x + b)\, T_0,
\]
et de l\`a
\[
a^2\,(a, b, c, d \Cdel y - b T_0, a T_0)^3 = 0
\]
ou enfin
\[
y^3 + 3 y\,(b^2 - a c)\, T_0{}^2 + (a^2 d - 3 a b c + 2 b^3)\, T_0{}^3 = 0.
\]
Les valeurs des coefficients peuvent \^etre compl\'et\'ees, eu \'egard \`a ce qu'ils doivent \^etre des invariants de $(a, b, c, d \Cdel \xi, \eta)^3$, $(T_0, T_1 \Cdel \eta, -\xi)$ (ou, comme on vient de le dire, des covariants de $(a, b, c, d \Cdel T_0, T_1)^3$). Mais ce qu'on peut r\'eduire la seconde \'equation des coefficients $T_0, T_1$ sont au nombre de deux, ce n'est que dans le cas particulier, o\`u il l'on aurait \`a une \'equation lin\'eaire.

\medskip

\hfill Londres, 18 Avril 1860.

% ============================================================
% PAPER 274
% DEUXIEME NOTE SUR LA TRANSFORMATION DE TSCHIRNHAUSEN
% ============================================================

\newpage
\noindent\textbf{368}\hfill\textbf{[274}

\medskip

\begin{center}
{\Large\bfseries 274.}\\[1em]
{\large\scshape Deuxi\`eme note sur la transformation de Tschirnhausen.}\\[1em]
{\small [From the \textit{Journal f\"ur die reine und angewandte Mathematik} (Crelle), tom. \textsc{lviii}.\ (1861), pp.~263--269.]}
\end{center}

\medskip

\noindent A la fin de ma premi\`ere note sur ce sujet j'ai appliqu\'e la transformation de Tschirnhausen \`a l'\'equation du troisi\`eme degr\'e mise sous la forme
\[
a x^3 + 3 b x^2 + 3 c x + d = (a, b, c, d \Cdel x, 1)^3 = 0.
\]
En y substituant $\tfrac{1}{3} b, \tfrac{1}{3} c$ au lieu de $b, c$, cette \'equation se change en
\[
a x^3 + b x^2 + c x + d = (a, b, c, d \Cdel x, 1)^3 = 0,
\]
et en m\^eme temps le r\'esultat obtenu dans ma premi\`ere note s'\'enonce de la mani\`ere suivante:

\medskip

En calculant pour l'\'equation
\[
(a, b, c, d \Cdel x, 1)^3 = 0,
\]
la transform\'ee en
\[
y = (a x + \tfrac{1}{3} b)\, T_0 + (a x^2 + b x + \tfrac{1}{3} c)\, T_1,
\]
on obtient
\[
\begin{aligned}
y^3 \quad &\\
{} + \tfrac{1}{3} y \quad \bigg\{ \quad &T_0{}^2\,(3 a c - b^2) \\
+\quad & T_0 T_1\,(9 a d - b c) \\
+\quad & T_1{}^2\,(3 b d - c^2) \bigg\} \\
{} + \tfrac{1}{27}\quad \bigg\{ \quad & T_0{}^3\,(27 a^2 d - 9 a b c + 2 b^3) \\
+\quad & T_0{}^2 T_1\,(27 a b d - 18 a c^2 + 3 b^2 c) \\
+\quad & T_0 T_1{}^2\,(-27 a c d + 18 b^2 d - 3 b c^2) \\
+\quad & T_1{}^3\,(-27 a d^2 + 9 b c d - 2 c^3) \bigg\} &= 0.
\end{aligned}
\]

\newpage
\noindent\textbf{274]}\hfill\textsc{deuxi\`eme note sur la transformation de tschirnhausen.}\hfill\textbf{369}

\medskip

\noindent Je vais me servir de cette formule, pour en d\'eduire l'\'equation qui d'une mani\`ere analogue est la transform\'ee de l'\'equation du quatri\`eme ordre
\[
(a, b, c, d, e \Cdel x, 1)^4 = 0.
\]
J'\'ecris d'abord
\[
(a, b, c, d, e \Cdel x, 1)^4 = 0,
\]
\[
y = (a x + \tfrac{1}{4} b)\, T_0 + (a x^2 + b x + \tfrac{1}{2} c)\, T_1 + (a x^3 + b x^2 + c x + \tfrac{1}{4} d)\, T_2,
\]
et je remarque qu'en faisant $e = 0$, le syst\`eme propos\'e se partage en deux, dont le premier est
\[
x = 0,
\]
\[
y = \tfrac{1}{4}(b T_0 + 2 c T_1 + 3 d T_2),
\]
et le second:
\[
(a, b, c, d \Cdel x, 1)^3 = 0,
\]
\[
y = (a x + \tfrac{1}{4} b)\, T_0 + (a x^2 + b x + \tfrac{1}{2} c)\, T_1 - \tfrac{1}{4} d\, T_2;
\]
ou, ce qui est la m\^eme chose,
\[
(a, b, c, d \Cdel x, 1)^3 = 0,
\]
\[
y + \tfrac{1}{4}(b T_0 + 2 c T_1 + 3 d T_2) = (a x + \tfrac{1}{4} b)\, T_0 + (a x^2 + b x + \tfrac{1}{2} c)\, T_1.
\]
Une circonstance analogue a lieu dans l'\'equation en $y$, r\'esultat de l'\'elimination du syst\`eme propos\'e. Pour $e = 0$ son premier membre se r\'esout de m\^eme en deux facteurs qui \'egal\'es \`a z\'ero sont les r\'esultats de l'\'elimination du premier et du second syst\`eme ci-dessus \'ecrits. Le premier de ces deux facteurs est donc
\[
y - \tfrac{1}{4}(b T_0 + 2 c T_1 + 3 d T_2),
\]
et le second (en vertu de la formule donn\'ee ant\'erieurement)
\[
\begin{aligned}
& \{y + \tfrac{1}{4}(b T_0 + 2 c T_1 + 3 d T_2)\}^3 \\
& \quad + \tfrac{1}{3}\big[(3 a c - b^2)\, T_0{}^2 + \text{etc.}\big]\,\{y + \tfrac{1}{4}(b T_0 + 2 c T_1 + 3 d T_2)\} \\
& \quad + \tfrac{1}{27}\big[(27 a^2 d - 9 a b c + 2 b^3)\, T_0{}^3 + \text{etc.}\big];
\end{aligned}
\]
donc en multipliant les deux facteurs, et en \'egalant \`a z\'ero leur produit, on a la transform\'ee en $y$ de la forme $(a, b, c, d, 0 \Cdel x, 1)^4$. Or dans le cas g\'en\'eral, o\`u $e$ est diff\'erent de z\'ero, les coefficients de la transform\'ee en $y$ sont des invariants des deux formes
\[
(a, b, c, d, e \Cdel \xi, \eta)^4,\quad (T_0, T_1, T_2 \Cdel \eta, -\xi)^2.
\]
Cette propri\'et\'e permet de d\'eduire leurs valeurs g\'en\'erales des valeurs particuli\`eres qu'ils ont pour $e = 0$. Je formerai de cette mani\`ere la transform\'ee en $y$ pour la forme $(a, b, c, d, 0 \Cdel x, 1)^4$, je passerai de l\`a \`a la forme $(a, b, c, d, 0 \Cdel x, 1)^4$ (ce qui se fait en \'ecrivant $4 b, 6 c, 4 d$ au lieu de $b, c, d$), et enfin je compl\'eterai les valeurs des coefficients en y introduisant $e$ au moyen de la propri\'et\'e que doivent poss\'eder les coefficients d'\^etre des invariants des deux formes
\[
(a, b, c, d, e \Cdel \xi, \eta)^4,\quad (T_0, T_1, T_2 \Cdel \eta, -\xi)^2.
\]
On obtient d'abord l'\'equation en $y$ sous la forme
\[
(1, 0, \mathfrak{C}, \mathfrak{D}, \mathfrak{E} \Cdel y, 1)^4 = 0,
\]

\hfill 47

\newpage
\noindent\textbf{370}\hfill\textsc{deuxi\`eme note sur la transformation de tschirnhausen.}\hfill\textbf{[274}

\medskip

\noindent o\`u
\[
\begin{aligned}
8\mathfrak{C} = {}& T_0{}^2\,(8 a c - 3 b^2) \\
& + T_0 T_1\,(24 a d - 4 b c) \\
& + T_1{}^2\,(8 b d - 4 c^2) \\
& + T_1 T_2\,(-2 b d) \\
& + T_2 T_3\,(-4 c d) \\[4pt]
8\mathfrak{D} = {}& T_0{}^3\,(8 a^2 d - 4 a b c + b^3) \\
& + T_0{}^2 T_1\,(4 a b d - 3 a c^2 + 2 b^2 c) \\
& + T_0{}^2 T_2\,(-2 a c d + 4 b^2 d - 4 a d^2) \\
& + T_0 T_1{}^2\,(-12 a d^2 + 16 b c^2) \\
& + T_0 T_1 T_2\,(-4 b d^2) \\
& + T_0 T_2{}^2\,(-7 b d^2) \\
& + T_1{}^2 T_2\,(-2 c d^2) \\
& + T_2{}^3\,(-1 d^3).
\end{aligned}
\]
\[
\begin{aligned}
256\mathfrak{E} = {}& T_0{}^4\,(64 a^2 b d + 16 a b^2 c - 3 b^4) \\
& + T_0{}^3 T_1\,(-128 a^2 c d - 80 a b^2 d + 64 a b c^2 - 8 b^3 c) \\
& + (T_0{}^3 T_2 \text{ and other powers, including:}) \\
& T_0{}^2 T_1{}^2\,(-128 a b c d + 64 a c^3 - 48 b^3 d + 8 b^2 c^2) \\
& + T_0{}^2 T_1 T_2\,(64 a b d^2 + 64 a c^2 d - 128 b c^2 d + 32 b c^3) \\
& + T_0{}^2 T_2{}^2\,(128 a c d^2 - 64 b c^2 d + 16 c^4) \\
& + T_0 T_1{}^3\,(-192 a^2 d^2 + 32 a b c d - 4 b^2 d^2) \\
& + T_0 T_1{}^2 T_2\,(-288 a b d^2 + 144 a c^2 d + 8 b^2 c d) \\
& + (\text{further terms}) \\
& - 160 b^2 d^2 + 48 b c^2 d \\
& + T_0 T_1 T_2{}^2\,(+192 a d^3 - 128 b c d^2 + 32 c^3 d) \\
& + T_0 T_2{}^3\,(-48 a c d^2 + 14 b^2 d^2) \\
& - 144 a c d^2 + 8 b c d^3 \\
& + T_1{}^4\,(-48 b d^3 + 8 c^3 d^2) \\
& + (\text{remaining terms}) \\
& + (-16 d^4) \\
& + T_2{}^4\,(-8 d^4).
\end{aligned}
\]

\noindent \textit{[Note: The coefficient table for $256\mathfrak{E}$ is densely packed in the original printing; the OCR layout shows several columns of monomial coefficients in $T_0,T_1,T_2$ which have been rendered above in approximate column-by-column form. The reader is referred to the scan for the precise tabular layout.]}

\medskip

Ce calcul achev\'e et substituant la forme $(a, b, c, d, 0 \Cdel x, 1)^4$ au lieu de $(a, b, c, d, 0 \Cdel x, 1)^4$ (o\`u $4 b, 6 c, 4 d$ au lieu de $b, c, d$) on obtient tous les termes de l'\'equation cherch\'ee, hormis ceux qui contiennent $e$: et ces derniers s'obtiennent au moyen de ce que les coefficients des diff\'erentes puissances de $y$ se r\'eduisent \`a z\'ero par l'op\'erateur
\[
a \partial_b + 2 b \partial_c + 3 c \partial_d + 4 d \partial_e - T_1 \partial_{T_0} - 2 T_2 \partial_{T_1}.
\]

\newpage
\noindent\textbf{274]}\hfill\textsc{deuxi\`eme note sur la transformation de tschirnhausen.}\hfill\textbf{371}

\medskip

\noindent Cela ne pr\'esente pas de difficult\'e, je supprime donc les calculs interm\'ediaires et je donne le r\'esultat final que voici: les \'equations
\[
(a, b, c, d, e \Cdel x, 1)^4 = 0,
\]
\[
y = (a x + b)\, T_0 + (a x^2 + 4 b x + 3 c)\, T_1 + (a x^3 + 4 b x^2 + 6 c x + 3 d)\, T_2,
\]
conduisent \`a la transform\'ee:
\[
(1, 0, \mathfrak{C}, \mathfrak{D}, \mathfrak{E} \Cdel y, 1)^4 = 0,
\]
o\`u l'on a (table dense en monomes $T_0^i T_1^j T_2^k$ avec coefficients en $a, b, c, d, e$):

\medskip

\textit{[Figure: large multi-column coefficient table for $\mathfrak{C}$, $\mathfrak{D}$, $\mathfrak{E}$ in monomes of $T_0, T_1, T_2$ with polynomial entries in $a, b, c, d, e$. The table spans most of pp.~371--372 of the original printing.]}

\medskip

\noindent The table for $\mathfrak{C}$ has the form (column entries are coefficients of $T_0{}^2, T_0 T_1, T_0 T_2, T_1{}^2, T_1 T_2, T_2{}^2$ respectively):
\[
\mathfrak{C} = 2 \begin{pmatrix}
a c + 8 \\
b^2 - 3
\end{pmatrix} T_0{}^2
+ \begin{pmatrix}
a d + 24 \\
b c - 8
\end{pmatrix} T_0 T_1
+ \begin{pmatrix}
a e + 32 \\
b d - 2 \\
c^2 - 6
\end{pmatrix} T_0 T_2
+ \begin{pmatrix}
a e + 16 \\
b d + 8 \\
c^2 - 6
\end{pmatrix} T_1{}^2
+ \begin{pmatrix}
b e + 24 \\
c d - 4
\end{pmatrix} T_1 T_2
+ \begin{pmatrix}
c e + 8 \\
d^2 - 3
\end{pmatrix} T_2{}^2.
\]

\medskip

\noindent Une table semblable s'applique \`a $\mathfrak{D}$ et \`a $\mathfrak{E}$, dont les expressions compl\`etes sont donn\'ees aux pp.~371--372 de l'\'edition originale. \textit{[Voir scan pour la disposition d\'etaill\'ee.]}

\medskip

J'\'ecris
\[
U' = a T_0{}^2 + 4 b T_0 T_1 + c\,(2 T_0 T_2 + 4 T_1{}^2) + 4 d T_1 T_2 + e T_2{}^2,
\]
\[
\begin{aligned}
H' = {}& (ac - b^2)\, T_0{}^2 + 2(ad - bc)\, T_0 T_1 + (ae - 2 bd + c^2)\, T_0 T_2 + 4(bd - c^2)\, T_1{}^2 \\
& + 2(be - cd)\, T_1 T_2 + (ce - d^2)\, T_2{}^2,
\end{aligned}
\]

\hfill $47$--$2$

\newpage
\noindent\textbf{372}\hfill\textsc{deuxi\`eme note sur la transformation de tschirnhausen.}\hfill\textbf{[274}

\medskip

\noindent et je repr\'esente par $4 \Phi'$ la valeur qui vient d'\^etre trouv\'ee pour $\mathfrak{D}$. Ces expressions $U'$, $H'$, $\Phi'$ sont des invariants des deux formes $(a, b, c, d, e \Cdel \xi, \eta)^4$, $(T_0, T_1, T_2 \Cdel \eta, -\xi)^2$.

\medskip

On a de plus les invariants
\[
ae - 4 bd + 3 c^2,\quad ace - ad^2 - b^2 e + 2 bcd - c^3,
\]
que je repr\'esente comme \`a l'ordinaire par $I, J$, et l'invariant $T_0 T_2 - T_1{}^2$ que je repr\'esente par $\Theta'$. Cela pos\'e on a
\[
\mathfrak{C} = 6 H' - 2 I \Theta',
\]
\[
\mathfrak{D} = 4 \Phi',
\]
\[
\mathfrak{E} = I U'^2 - 3 H'^2 + I^2 \Theta'^2 + 12 J \Theta' U' + 2 I \Theta' H'.
\]
La derni\`ere de ces \'equations peut \^etre v\'erifi\'ee ais\'ement, pour cela on a seulement besoin de remarquer qu'en posant $a = e = 1$, $b = d = 0$, $c = \theta$, elle devient
\[
\begin{aligned}
&(1 + 3\theta^2)\,\{T_0{}^2 + \theta\,(2 T_0 T_2 + 4 T_1{}^2) + T_2{}^2\}^2 \\
&\quad - 3\,\{\theta T_0{}^2 + (1 + \theta^2)\, T_0 T_2 - 4 \theta^2 T_1{}^2 + \theta T_2{}^2\}^2 \\
&\quad + (1 + 3\theta^2)^2\,(T_0 T_2 - T_1{}^2)^2 \\
&\quad + 12\,(\theta - \theta^3)\,(T_0 T_2 - T_1{}^2)\,(T_0{}^2 + \theta\,(2 T_0 T_2 + 4 T_1{}^2) + T_2{}^2) \\
&\quad + 2\,(1 + 3\theta^2)\,(T_0 T_2 - T_1{}^2)\,(\theta T_0{}^2 + (1 + \theta^2)\, T_0 T_2 - 4 \theta^2 T_1{}^2 + \theta T_2{}^2)
\end{aligned}
\]
\[
\begin{aligned}
={}& T_0{}^4 \\
& + T_0{}^3 T_1\,(12 \theta) \\
& + T_0{}^2 T_1{}^2\,(-6 \theta^2 + 54 \theta^4) \\
& + T_0{}^2 T_2{}^2\,(2 + 36 \theta^2) \\
& + T_0 T_1{}^2 T_2\,(-4 + 36 \theta^2) \\
& + T_1{}^4\,(-18 \theta^2 + 81 \theta^4) \\
& + T_2{}^2 T_1{}^2\,(-6 \theta^2 + 54 \theta^4) \\
& + T_0 T_2{}^3\,(12 \theta)
\end{aligned}
\]
\'equation qui est identique. L'expression de l'invariant $I$ (quadrinvariant) de la fonction $(1, 0, \mathfrak{C}, \mathfrak{D}, \mathfrak{E} \Cdel y, 1)^4$ est $\mathfrak{E} + 3 \mathfrak{C}^2$ ou $\mathfrak{E} + 3(H' - \tfrac{1}{3} I \Theta')^2$, c'est-\`a-dire
\[
I U'^2 - 3 H'^2 + I^2 \Theta'^2 + 12 J \Theta' U' + 2 I \Theta' H' + 3 H'^2 + \tfrac{1}{3} I^2 \Theta'^2 - 2 I \Theta' H',
\]
ou enfin
\[
I U'^2 + \tfrac{4}{3} I^2 \Theta'^2 + 12 J \Theta' U',
\]

\newpage
\noindent\textbf{274]}\hfill\textsc{deuxi\`eme note sur la transformation de tschirnhausen.}\hfill\textbf{373}

\medskip

\noindent ce qui est \'egal \`a
\[
\tfrac{1}{I}\,\big[(I U' + 6 J \Theta')^2 + \tfrac{4}{3}(I^3 - 27 J^2)\, \Theta'^2\big].
\]
La condition \`a remplir pour que cet invariant se r\'eduise \`a z\'ero peut donc \^etre pr\'esent\'ee sous la forme
\[
I U' + [6 J \pm 2 \sqrt{-\tfrac{1}{3}(I^3 - 27 J^2)}]\, \Theta' = 0,
\]
ce qui s'accorde avec un r\'esultat trouv\'e par M.~Hermite.

\medskip

Il doit y avoir, ce me semble, une \'equation identique de la forme
\[
J U'^3 - I U'^2 H' + I^2 \Theta' U' + M U' = -\Phi'^2
\]
qui servirait \`a exprimer le carr\'e de $\Phi'$ au moyen des autres invariants $U'$, $H'$, $\Theta'$, $I$, $J$, mais en supposant que cette \'equation existe, la forme du facteur $M$, que je n'ai pas encore cherch\'ee, reste \`a d\'eterminer; l'invariant $J$ (cubinvariant) de la forme $(1, 0, \mathfrak{C}, \mathfrak{D}, \mathfrak{E} \Cdel y, 1)^4$ contient $\Phi'^2$ et il faudrait employer l'identit\'e dont je viens de parler pour r\'eduire \`a sa forme la plus simple cet invariant; dans l'\'etat actuel de la question je ne m'occupe donc pas de l'expression du cubinvariant de $(1, 0, \mathfrak{C}, \mathfrak{D}, \mathfrak{E} \Cdel y, 1)^4$.

\medskip

Pour passer au cas d'une \'equation du cinqui\`eme ordre, on devra faire usage de la formule qui se rapporte \`a la forme $(a, b, c, d, e \Cdel x, 1)^4$. En faisant la substitution n\'ecessaire on arrive \`a ce r\'esultat que pour l'\'equation
\[
(a, b, c, d, e \Cdel x, 1)^4 = 0
\]
la transform\'ee en
\[
y = (a x + \tfrac{1}{4} b)\, T_0 + (a x^2 + b x + \tfrac{1}{2} c)\, T_1 + (a x^3 + b x^2 + c x + \tfrac{1}{4} d)\, T_2
\]
est la suivante
\[
(1, 0, \mathfrak{C}, \mathfrak{D}, \mathfrak{E} \Cdel y, 1)^4 = 0.
\]

\medskip

\textit{[Figure: large multi-column tables giving $\mathfrak{C}$, $\mathfrak{D}$, $\mathfrak{E}$ as explicit polynomials in $T_0, T_1, T_2$ with coefficients formed from monomes of $a, b, c, d, e$ such as $ac+8$, $b^2-3$, $a^2bd+8$, $abce-60$, etc. The tables are densely printed on pp.~373--374 of the original and contain hundreds of monomial entries grouped column-by-column under each $T_0{}^i T_1{}^j T_2{}^k$ heading. The reader is referred to the original scan for the full layout.]}

\medskip

The leading entries (samples from the columns) are:
\[
\mathfrak{C} = \tfrac{1}{4}\!\left[
\begin{array}{lll}
T_0{}^2: & ac + 8,\ b^2 - 3; \\
T_0 T_1: & ad + 24,\ bc - 4; \\
T_0 T_2: & ae + 32,\ bd - 2,\ c^2 - 6; \\
T_1{}^2: & ae + 16,\ bd + 8,\ c^2 - 6; \\
T_1 T_2: & be + 24,\ cd - 4; \\
T_2{}^2: & ce + 8,\ d^2 - 3.
\end{array}
\right]
\]

\medskip

\textit{[Figure: similar columnar tables for $\mathfrak{D}$ and $\mathfrak{E}$ with respective leading entries $T_0{}^3: a^2 d + 8$, $abc - 4$, $b^3 + 1$, etc., and $T_0{}^4: a^3 e + 256$, $a^2 bd - 64$, $a b^2 c + 16$, $b^4 - 3$, etc. The complete tables are reproduced in the scanned original.]}

\newpage
\noindent\textbf{374}\hfill\textsc{deuxi\`eme note sur la transformation de tschirnhausen.}\hfill\textbf{[274}

\medskip

\noindent \textit{[Figure: continuation of the multi-column coefficient tables for $\mathfrak{E}$. Each column carries a header of the form $T_0{}^i T_1{}^j T_2{}^k$ (with $i+j+k=4$), and the column entries are integer-weighted monomes of $a, b, c, d, e$ summing to the coefficient of that monome in $256 \mathfrak{E}$. Sample headers include $T_0{}^4$, $T_0{}^3 T_1$, $T_0{}^3 T_2$, $T_0{}^2 T_1{}^2$, $T_0{}^2 T_1 T_2$, $T_0{}^2 T_2{}^2$, $T_0 T_1{}^3$, $T_0 T_1{}^2 T_2$, $T_0 T_1 T_2{}^2$, $T_0 T_2{}^3$, $T_1{}^4$, $T_1{}^3 T_2$, $T_1{}^2 T_2{}^2$, $T_1 T_2{}^3$, $T_2{}^4$.]}

\medskip

En m'appuyant sur ce r\'esultat j'esp\`ere \^etre \`a m\^eme de trouver la formule pour l'\'equation du cinqui\`eme ordre.

\medskip

\hfill Londres, 11 Mai 1860.

% ============================================================
% PAPER 275
% ON TSCHIRNHAUSEN'S TRANSFORMATION
% ============================================================

\newpage
\noindent\textbf{275]}\hfill\textbf{375}

\medskip

\begin{center}
{\Large\bfseries 275.}\\[1em]
{\large\scshape On Tschirnhausen's Transformation.}\\[1em]
{\small [From the \textit{Philosophical Transactions of the Royal Society of London}, vol. \textsc{cli}.\ (for the year 1861), pp.~561--578. Received November 7,---Read December 5, 1861.]}
\end{center}

\medskip

\noindent The memoir of M.~Hermite, ``Sur quelques th\'eor\`emes d'alg\`ebre et la r\'esolution de l'\'equation du quatri\`eme degr\'e,'' \textit{Comptes Rendus}, t.~\textsc{xlvi}.\ p.~961 (1858), contains a very important theorem in relation to Tschirnhausen's Transformation of an equation $f(x) = 0$ into another of the same degree in $y$, by means of the substitution $y = \phi x$, where $\phi$ is a rational and integral function of $x$. In fact, considering for greater simplicity a quartic equation,
\[
(a, b, c, d, e \Cdel x, 1)^4 = 0,
\]
M.~Hermite gives to the equation $y = \phi x$ the following form,
\[
y = a T + (a x + 4 b)\, B + (a x^2 + 4 b x + 6 c)\, C + (a x^3 + 4 b x^2 + 6 c x + 4 d)\, D,
\]
(I write $B$, $C$, $D$ in the place of his $T_0$, $T_1$, $T_2$), and he shows that the transformed equation in $y$ has the following property: viz., every function of the coefficients which, expressed as a function of $a, b, c, d, e, T, B, C, D$, does not contain $T$, is an invariant, that is, an invariant of the two quantics
\[
(a, b, c, d, e \Cdel X, Y)^4,\quad (B, C, D \Cdel Y, -X)^2.
\]
This comes to saying that if $T$ be so determined that in the equation for $y$ the coefficient of the second term ($y^3$) shall vanish, the other coefficients will be invariants; or if, in the function of $y$ which is equated to zero, we consider $y$ as an absolute constant, the function of $y$ will be an invariant of the two quantics. It is easy to find the value of $T$; this is in fact given by the equation
\[
0 = a T + 3 b B + 3 c C + d D;
\]

\newpage
\noindent\textbf{376}\hfill\textsc{on tschirnhausen's transformation.}\hfill\textbf{[275}

\medskip

\noindent and we have thence for the value of $y$,
\[
y = (a x + b)\, B + (a x^2 + 4 b x + 3 c)\, C + (a x^3 + 4 b x^2 + 6 c x + 3 d)\, D;
\]
so that for this value of $y$ the function of $y$ which equated to zero gives the transformed equation will be an invariant of the two quantics. It is proper to notice that in the last-mentioned expression for $y$, all the coefficients except those of the term in $x^0$, or $b B + 3 c C + 3 d D$, are those of the binomial $(1, 1)^4$, whereas the excepted coefficients are those of the binomial $(1, 1)^3$; this suffices to show what the expression for $y$ is in the general case.

\medskip

I have in the two papers, ``Note sur la Transformation de Tschirnhausen,'' and ``Deuxi\`eme Note sur la Transformation de Tschirnhausen,'' Crelle, t.~\textsc{lviii}.\ pp.~259 and 263 (1861), [273 and 274], obtained the transformed equations for the cubic and quartic equations; and by means of a grant from the Government Grant Fund, I have been enabled to procure the calculation, by Messrs Davis and Otter, under my superintendence, of the transformed equation for the quintic equation. The several results are given in the present memoir; and for greater completeness, I reproduce the demonstration which I have given in the former of the above-mentioned two Notes, of the general property, that the function of $y$ is an invariant. At the end of the memoir I consider the problem of the reduction of the general quintic equation to Mr~Jerrard's form $x^5 + a x + b = 0$.

\medskip

Considering for simplicity the foregoing two equations
\[
(a, b, c, d, e \Cdel x, 1)^4 = 0,
\]
\[
y = (a x + b)\, B + (a x^2 + 4 b x + 3 c)\, C + (a x^3 + 4 b x^2 + 6 c x + 3 d)\, D;
\]
let the second of these be represented by $y = V$, the transformed equation in $y$ is
\[
(y - V_1)(y - V_2)(y - V_3)(y - V_4) = 0,
\]
where $V_1, V_2, V_3, V_4$ are what $V$ becomes upon substituting therein for $x$ the roots $x_1, x_2, x_3, x_4$ of the quartic equation respectively. Considering $y$ as a constant, the conditions to be satisfied in order that the function in $y$ may be an invariant are that this function shall be reduced to zero by each of the two operators
\[
a \partial_b + 2 b \partial_c + 3 c \partial_d + 4 d \partial_e - (B \partial_C + 2 C \partial_D),
\]
\[
4 b \partial_a + 3 c \partial_b + 2 d \partial_c + e \partial_d - (2 C \partial_B + B \partial_C).
\]
These conditions will be satisfied if each of the factors $y - V_1$, \&c.\ has the property in question; that is, if $y - V$, or (what is the same thing) if $V$, supposing that $x$ denotes therein a root of the quartic equation, is reduced to zero by each of the two operators. Considering the first operator, which for shortness I represent by
\[
\Delta - (D \partial_C + 2 C \partial_B),
\]
in order to obtain $\Delta V$ we require the value of $\Delta x$. To find it, operating with $\Delta$ on the quartic equation, we have
\[
(a, b, c, d \Cdel x, 1)^3\, \Delta x + (a, b, c, d \Cdel x, 1)^3 = 0,
\]

\newpage
\noindent\textbf{275]}\hfill\textsc{on tschirnhausen's transformation.}\hfill\textbf{377}

\medskip

\noindent or $\Delta x = -1$. In $\Delta V$, the part which depends on the variation of $\Delta x$ then is
\[
-a B + (-2 a x - 4 b)\, C + (-3 a x^2 - 8 b x - 6 c)\, D,
\]
and the other part of $\Delta V$ is at once found to be
\[
+ a B + (4 a x + 6 b)\, C + (4 a x^2 + 12 b x + 9 c)\, D;
\]
whence, adding,
\[
\Delta V = 2(a x + b)\, C + (a x^2 + 4 b x + 3 c)\, D,
\]
and this is precisely equal to
\[
(D \partial_C + 2 C \partial_B)\, V;
\]
so that $V$ is reduced to zero by the operator $\Delta - (D \partial_C + 2 C \partial_B)$.

\medskip

Similarly, if the second operator is represented by
\[
\nabla - (2 C \partial_D + B \partial_C),
\]
then we have
\[
(a, b, c, d \Cdel x, 1)^3\, \nabla x + x\,(b, c, d, e \Cdel x, 1)^3 = 0,
\]
which by means of the equation
\[
(a, b, c, d, e \Cdel x, 1)^4 = 0
\]
is reduced to $\nabla x = x^2$. Hence in $\nabla V$ the part depending on the variation of $x$ is
\[
(a x^2)\, B + (2 a x^3 + 4 b x^2)\, C + (3 a x^4 + 8 b x^3 + 6 c x^2)\, D,
\]
and the other part of $\nabla V$ is at once found to be
\[
(4 b x + 3 c)\, B + (4 b x^2 + 12 c x + 6 d)\, C + (4 b x^3 + 12 c x^2 + 12 d x + 3 e)\, D;
\]
and, adding, the coefficient of $D$ vanishes on account of the quartic equation, and we have
\[
\nabla V = (a x^2 + 4 b x + 3 c)\, B + 2(a x^3 + 4 b x^2 + 6 c x + 3 d)\, C,
\]
which is precisely equal to
\[
(2 C \partial_D + B \partial_C)\, V,
\]
so that $V$ is reduced to zero by the operator
\[
\nabla - (2 C \partial_D + B \partial_C),
\]
which completes the demonstration; and the demonstration in the general case is precisely similar.

\medskip

In the case of the cubic equation we have
\[
(a, b, c, d \Cdel x, 1)^3 = 0,
\]
\[
y = (a x + b)\, B + (a x^2 + 3 b x + 2 c)\, C;
\]

\hfill C, IV.\hfill 48

\newpage
\noindent\textbf{378}\hfill\textsc{on tschirnhausen's transformation.}\hfill\textbf{[275}

\medskip

\noindent and writing the second equation in the form
\[
(y - b B - 2 c C) + x\,(-a B - 3 b C) + x^2\,(-a C) = 0,
\]
multiplying by $x$ and reducing by the cubic equation, we have
\[
dC + x\,(y - b B + c C) + x^2\,(-a B) = 0,
\]
and repeating the process,
\[
dB + x\,(3 c B + d C) + x^2\,(y + 2 b B + c C) = 0;
\]
or, what is the same thing, we have the system of equations
\[
\begin{vmatrix}
y - b B - 2 c C, & -a B - 3 b C, & -a C \\
d C, & y - b B + c C, & -a B \\
d B, & 3 c B + d C, & y + 2 b B + c C
\end{vmatrix}\!(1, x, x^2) = 0,
\]
and the resulting equation in $y$ is of course that formed by equating to zero the determinant formed out of the matrix in this equation. The developed expression is
\[
(1, 0, \mathfrak{C}, \mathfrak{D} \Cdel y, 1)^3 = 0,
\]
where
The coefficients $\mathfrak{C}$ and $\mathfrak{D}$ are given column-wise in the original printing as follows:
\[
\mathfrak{C} = \tfrac{1}{3}\big[\, (ac - b^2)\, B^2 + (ad - 3 bc)\, BC + (-3 bd + c^2)\, C^2\,\big] \quad (\text{column sums } -2,\ -3,\ -2),
\]
\[
\mathfrak{D} = \tfrac{1}{27}\big[\, (a^2 d - 3 abc + 2 b^3)\, B^3 + (3 abd + 3 b^2 c)\, B^2 C + (6 b^2 d + 6 acd)\, B C^2 + (a d^2 + 3 bcd + 3)\, C^3\,\big].
\]

\medskip

\textit{[Figure: Cayley's column-table for the coefficients $\mathfrak{C}$ and $\mathfrak{D}$. The columns are headed $B^2, BC, C^2$ (for $\mathfrak{C}$) and $B^3, B^2 C, B C^2, C^3$ (for $\mathfrak{D}$); within each column the entries are monomes in $a, b, c, d$ with their numerical coefficients, the column sums being zero as a check (the original adds annotations like $+1, -1, +3$ etc., giving the positive and negative subtotals for each column).]}

\medskip

The sum of the coefficients in each column should here and elsewhere in the present memoir be equal to zero, and I have by way of verification annexed to each column the sums ($\pm$ a number) of the positive and negative coefficients. The coefficients $\mathfrak{C}, \mathfrak{D}$, and therefore the function in $y$, are invariants of the two forms,
\[
(a, b, c, d \Cdel X, Y)^3,\quad (B, C \Cdel Y, -X)^2;
\]
or in the present case, where there are only two coefficients $B, C$, the coefficients $\mathfrak{C}, \mathfrak{D}$, and therefore also the function in $y$, are covariants of the single form $(a, b, c, d \Cdel B, C)^3$, considering therein $(B, C)$ as the facients.

\end{document}
